%==============================================================================
% Experiencia Profesional
%==============================================================================
\section{Experiencia Profesional}

\cventry{Agosto 2026--Presente}
        {Analista Funcional y Desarrollador de Software}
        {COTECMAR}
        {Cartagena, Colombia}{}
{Corporación de Ciencia y Tecnología para el Desarrollo de la Industria Naval.
\begin{itemize}
  \item Análisis funcional y diseño de la arquitectura de \textbf{TOP} (Portal de
        Operaciones Tecnológicas), sistema tipo MES que opera entre el piso de
        producción del astillero y el ERP corporativo SAP, en reemplazo de una
        gestión fragmentada en archivos Excel y cronogramas aislados.
  \item Ingeniería inversa del sistema \emph{legacy} y rediseño del modelo relacional
        (ERD documentado en Mermaid.js): unificación de estructuras redundantes de
        costos en una jerarquía única y corrección del modelo para soportar la
        realidad comercial de un proyecto con múltiples cascos.
  \item Diseño del modelo de integración con \textbf{SAP MM/PS}: inyección de
        solicitudes de pedido y reservas desde TOP, con una regla de caducidad que
        elimina las reservas fantasma responsables de inmovilizar inventario.
  \item Definición del indicador operativo de productividad (\textbf{Kg/HH}) como
        valor derivado, calculado sobre horas-hombre efectivamente auditadas y
        aprobadas, y del circuito de validación Supervisor $\rightarrow$ Jefe de Área
        previo a la imputación de consumos.
  \item Rediseño funcional del módulo comercial de estimaciones: máquina de ocho
        estados con la firma como prerrequisito del contrato, y vistas Kanban,
        tabla y calendario.
  \item Desarrollo sobre \textbf{Laravel 10} (PHP 8), \textbf{Vue 3} e
        \textbf{Inertia.js}, con MySQL como base operativa y conexiones SQL Server
        hacia SAP.
\end{itemize}}

\cventry{Agosto 2025--Mayo 2026}
        {Desarrollador Full-Stack --- Proyecto de Ingeniería de Sistemas}
        {VIT Cloud Services}
        {Bogotá, Colombia}{}
{\begin{itemize}
  \item Desarrollo full-stack de una plataforma web para la gestión de la
        Seguridad de la Información corporativa (SGSI), alineada a la norma
        \textbf{ISO/IEC 27001:2022}.
  \item Lideré la definición de requerimientos funcionales en comunicación
        directa con la gerencia del cliente, traduciendo los controles de la
        norma en flujos de trabajo y modelo de datos.
  \item Construcción con \textbf{Django}, \textbf{PostgreSQL} y \textbf{React}
        a lo largo de dos semestres académicos.
\end{itemize}}

% --- Duplicar este bloque para experiencias anteriores, prácticas o freelance ---
% \cventry{[AAAA]--[AAAA]}{[Cargo]}{[Empresa]}{[Ciudad]}{}{[Logros medibles]}
